\section{Power decay of the selected charge energy}
\label{sec:arealaw-source-selected-energy-power}

Fix the final scanner exponents
$\omega=1/20000$, $\ell=1/100000$, $\kappa=1/1000000$
and $\nu=1/1000$ of \cite{OpenAI2026AreaLaw}.
The charge-energy expression in Section~9, lines~492--511, contains
the eighth root of a normalized entropy defect. Thus a density bound
at rate $8\nu$ suffices to put that term at rate $\nu$.

\begin{definition}[Final-scale charge energy bound]
    \label{def:arealaw-final-selected-energy-cap}
    \lean{TNLean.PEPS.AreaLaw.Scan.finalSelectedEnergyCap}
    For $C,C_\ell,P,W,x,p\in\mathbb R$ and an integer $n\geq1$, set
    $D=\lceil n^\kappa\rceil$ and
    \begin{align}
        E_n(W,x,p)
         & =CW^2n^\ell D^2(\log n)^{C_\ell}
        \left((1-p)x^{1/8}+(1-p)n^{-P}+p\right).
        \label{eq:selected_energy_cap}
    \end{align}
    Here $x$ denotes the normalized entropy defect.
\end{definition}

\begin{theorem}[Exact energy exponents]
    \label{thm:arealaw-selected-energy-margins}
    \lean{TNLean.PEPS.AreaLaw.Scan.final_selectedEnergy_margins}
    At the final scanner exponents,
    \begin{align}
        2\omega+\ell+2\kappa     & =\frac{14}{125000},   &
        2\omega+\ell+2\kappa-\nu & =-\frac{111}{125000}, &
        -\frac1{2000}-(2\omega+\ell+2\kappa-\nu)
                                 & =\frac{97}{250000}.
        \label{eq:selected_energy_margins}
    \end{align}
\end{theorem}

\begin{proof}
    Substitute the displayed rational values and simplify.
\end{proof}

\begin{theorem}[Pointwise charge energy bound]
    \label{thm:arealaw-selected-energy-pointwise}
    \lean{TNLean.PEPS.AreaLaw.Scan.finalSelectedEnergyCap_le}
    Fix $C\geq0$, $C_\ell\in\mathbb R$ and $P\geq\nu$.
    For every integer $n\geq1$, every $1\leq W\leq n^\omega$,
    every $0\leq x\leq n^{-8\nu}$ and every
    $p\in[n^{-\nu},2n^{-\nu}]$, the expression in
    (\ref{eq:selected_energy_cap}) satisfies
    \begin{align}
        E_n(W,x,p) & \leq16C n^{-111/125000}(\log n)^{C_\ell}.
        \label{eq:selected_energy_pointwise}
    \end{align}
\end{theorem}

\begin{proof}
    The density bound gives $x^{1/8}\leq n^{-\nu}$, and
    $P\geq\nu$ gives $n^{-P}\leq n^{-\nu}$. Since $p\geq0$,
    multiplication by $1-p$ cannot increase either nonnegative term.
    Consequently
    \begin{align}
        (1-p)x^{1/8}+(1-p)n^{-P}+p & \leq4n^{-\nu}.
        \label{eq:selected_energy_bracket}
    \end{align}
    Also $D\leq n^\kappa+1\leq2n^\kappa$ and
    $W^2\leq n^{2\omega}$. Substitute these inequalities into
    (\ref{eq:selected_energy_cap}) and add the exponents to obtain
    (\ref{eq:selected_energy_pointwise}).
\end{proof}

\begin{theorem}[Uniform normalized energy decay]
    \label{thm:arealaw-selected-energy-decay}
    \lean{TNLean.PEPS.AreaLaw.Scan.eventually_final_selectedEnergy_le_rpow,TNLean.PEPS.AreaLaw.Scan.exists_final_selectedEnergy_le_rpow}
    Fix $C\geq0$, $C_\ell\in\mathbb R$, $g>0$ and $P\geq\nu$.
    There exists an integer $N\geq2$ such that, for every $n\geq N$,
    every $1\leq W\leq n^\omega$, every $0\leq x\leq n^{-8\nu}$,
    every $p\in[n^{-\nu},2n^{-\nu}]$ and every $e\leq n^{-1000}$,
    \begin{align}
        \frac{e+E_n(W,x,p)}{g/2} & \leq n^{-1/2000}.
        \label{eq:selected_energy_decay}
    \end{align}
    The threshold is chosen before $n,W,x,p,e$.
\end{theorem}

\begin{proof}
    \uses{thm:arealaw-selected-energy-pointwise}
    The strict gap $97/250000>0$ in
    (\ref{eq:selected_energy_margins}) absorbs every fixed coefficient
    and logarithmic power. Thus, eventually,
    \begin{align}
        16C n^{-111/125000}(\log n)^{C_\ell}
                  & \leq\frac g4 n^{-1/2000}, &
        n^{-1000} & \leq\frac g4 n^{-1/2000}.
        \label{eq:selected_energy_absorption}
    \end{align}
    Both estimates depend only on the fixed data. Apply
    (\ref{eq:selected_energy_pointwise}), add the inequalities in
    (\ref{eq:selected_energy_absorption}), and divide by $g/2>0$.
\end{proof}

\begin{theorem}[Separate replica remainder]
    \label{thm:arealaw-selected-energy-remainder}
    \lean{TNLean.PEPS.AreaLaw.Scan.selectedEnergy_add_remainder_le,TNLean.PEPS.AreaLaw.Scan.exists_final_selectedEnergy_add_remainder_le_rpow}
    Under the hypotheses of Theorem~\ref{thm:arealaw-selected-energy-decay},
    the same threshold works for every additional real remainder $r$:
    \begin{align}
        \frac{e+E_n(W,x,p)+r}{g/2}
         & \leq n^{-1/2000}+\frac r{g/2}.
        \label{eq:selected_energy_remainder}
    \end{align}
\end{theorem}

\begin{proof}
    \uses{thm:arealaw-selected-energy-decay}
    Expand the quotient of the sum and add $r/(g/2)$ to
    (\ref{eq:selected_energy_decay}).
\end{proof}

These conclusions are numerical. Applying them to a physical scanner
requires the normalized density and energy comparison at the same selected
charge point. If $r=r_k$ tends to zero as the replica count $k$ tends to
infinity, (\ref{eq:selected_energy_remainder}) keeps that limit separate
from the power estimate in $n$.
